%==============================================================================
\chapter{Objectifs du projet}
%==============================================================================

\section{Objectif principal}
Fournir une application \textbf{multi-niveaux, évolutive et faiblement couplée}
permettant de gérer l'ensemble du cycle de vie des ruches, de leur composition
physique jusqu'au suivi de production, et d'offrir des tableaux de bord d'aide
à la décision aux différents profils d'utilisateurs.

\section{Objectifs spécifiques}
\begin{itemize}
  \item Assurer les opérations \textbf{CRUD} (création, lecture, mise à jour, suppression) sur l'ensemble des entités métier.
  \item Gérer la hiérarchie \emph{fermier / ferme / site / ruche / corps ou hausse / cadre}.
  \item Planifier, approuver et suivre les visites d'inspection des agents apiculteurs.
  \item Produire des \textbf{rapports de visite} structurés et exploitables.
  \item Offrir des \textbf{tableaux de bord} filtrables (calendrier agents × ruches, alertes de production, alertes sanitaires, synthèse financière).
  \item Anticiper l'intégration de mesures \textbf{horodatées et géolocalisées} issues de capteurs à unités hétérogènes.
  \item Garantir \textbf{internationalisation}, \textbf{sécurité} et \textbf{interopérabilité} (services web).
\end{itemize}

\section{Résultats attendus (indicateurs de réussite)}
\begin{itemize}
  \item Support d'au moins \textbf{trois langues} (FR, EN, AR) en version de démonstration.
  \item Détection automatique des dépassements de seuils \textbf{paramétrables} (production, effectif).
  \item Restitution des données de production par ferme, site et ruche sous forme de graphiques.
  \item Exposition d'une API interrogeable par des applications tierces (Excel, etc.).
\end{itemize}

